% Generated by scripts/24_program_overview.py; do not edit.
\begin{longtable}{@{}>{\raggedright\arraybackslash}p{2.5cm}>{\raggedright\arraybackslash}p{1.9cm}>{\raggedright\arraybackslash}p{3.4cm}>{\raggedright\arraybackslash}p{3.2cm}>{\raggedright\arraybackslash}p{3.6cm}>{\raggedright\arraybackslash}p{6.6cm}>{\raggedright\arraybackslash}p{2.6cm}@{}}
\caption{Curated programs: source (ZFIN anatomy terms grouped by \protect\nolinkurl{data/curated/zfin_anatomy_programs.tsv}, or the signalling and state tables), the stages at which ZFIN records it (anatomy programs; all curated programs are defined at every stage), the program of Table~\ref{tab:programs} it corresponds to and the programs it is presumed to derive from (the arrows of Figure~\ref{fig:program-timeline}), where and under which signals it arises, its role for the model and interpretation, and its first genes (by mean heterogeneity over the stages, then records). Annotation: \protect\nolinkurl{data/curated/curated_programs.tsv}.}\label{tab:curated-programs}\\
\toprule
Program & ZFIN stages & Source & Table~\ref{tab:programs} / from & Location, signals & Role and interpretation & First genes \\
\midrule
\endfirsthead
\toprule
Program & ZFIN stages & Source & Table~\ref{tab:programs} / from & Location, signals & Role and interpretation & First genes \\
\midrule
\endhead
\bottomrule
\endlastfoot
\multicolumn{7}{@{}l}{\textit{early states}} \\
maternal &  & literature table & Maternal & whole embryo; degraded after MBT & Library-level timing control. Maternal transcripts; acts on whole libraries~\cite{white2017,wibrand2002,helde1993}. & \emph{h1m, cth1, buc, gdf3, mos, cldnd} \\
\rowcolor{gray!10}blastomere &  & literature table & Blastomeres; from maternal & whole blastoderm after MBT & The undifferentiated state from which the first decisions start. Pluripotency genes (pou5f3, nanog, soxB1); lost from 60\% epiboly in most cells~\cite{lee2013,hauptmann1995,xu2012,okuda2006}. & \emph{pou5f3, nanog, sox19b, lin28a, sox3, sox19a} \\
\midrule
\multicolumn{7}{@{}l}{\textit{blastoderm and margin}} \\
EVL & dome -- 6-somite & ZFIN: EVL; periderm; superficial blastomere & EVL; from blastomere & outermost cell layer of the blastoderm & First lineage split after MBT; a distinct population at every stage. Enveloping layer, an epithelium covering the embryo; becomes the periderm; a clear switch-like program~\cite{sagerstrom2005,kotkamp2014,sabel2009,siddiqui2010}. & \emph{krt4, krt18, cldne, cldnb, lye, krt8} \\
\rowcolor{gray!10}PGC & 30\% epiboly -- 6-somite & ZFIN: primordial germ cell; germ line cell & Primordial germ cells; from maternal & scattered clusters; specified by inherited germ plasm & A population outside the somatic decisions; a positive control for a rare population. Primordial germ cells; few cells per embryo~\cite{koprunner2001,weidinger2003,hashimoto2004}. & \emph{nanos3, ddx4, dnd1, ca15b, tdrd7a, gra} \\
margin & dome -- bud & ZFIN: margin; germ ring; marginal blastomere & Marginal mesendoderm; from blastomere & vegetal edge of the blastoderm (germ ring); Nodal from the YSL and the margin itself & The mesendoderm before internalization; the hypoblast/epiblast decision around 6 hpf starts here. Marginal cells that internalize to form mesoderm and endoderm; split by fate along the margin~\cite{rodaway1999,feldman1998,hagos2007}. & \emph{gsc, ved, ta, aplnra, mespab, eve1} \\
\rowcolor{gray!10}dorsal organizer & 30\% epiboly -- 75\% epiboly & ZFIN: shield; presumptive shield; presumptive dorsal mesoderm & Dorsal organizer; from margin & dorsal margin; maternal beta-catenin targets, then Nodal & The dorsal-ventral decision: the organizer protects the dorsal side from BMP. The shield (Spemann organizer); gives rise to prechordal plate and notochord~\cite{yamanaka1998,shimizu2002,stachel1993,melby2000}. & \emph{gsc, noto, insb, xbp1, pitx2, twist2} \\
dorsal forerunner & 50\% epiboly -- 6-somite & ZFIN: forerunner cell group; Kupffer's vesicle & Dorsal forerunner cells; from EVL & dorsal margin, ahead of the shield & A small population at the focus stages (shield, 60\% epiboly). Dorsal forerunner cells, which form the Kupffer's vesicle (left-right organ); shares genes with EVL and endoderm~\cite{melby1996}. & \emph{sult6b1, cxcr4a, rassf7b, vgll4l, foxj1a, si:ch73-234b20.5} \\
\midrule
\multicolumn{7}{@{}l}{\textit{endoderm and mesoderm}} \\
endoderm & 30\% epiboly -- 6-somite & ZFIN: endoderm; presumptive endoderm; endodermal cell & Endoderm; from margin & margin, closest to the edge; highest Nodal & The endoderm/mesoderm decision of the hypoblast. Endoderm (sox32/sox17 cells); gut and its derivatives~\cite{alexander1999,kikuchi2000,kikuchi2001,mizoguchi2008}. & \emph{gata5, foxa2, ackr3b, sox32, mixl1, sox17} \\
\rowcolor{gray!10}mesoderm & 30\% epiboly -- 6-somite & ZFIN: hypoblast; mesoderm; presumptive mesoderm & Marginal mesendoderm (internalized); from margin & hypoblast; Nodal, FGF, Wnt and BMP along the margin & The internalized side of the hypoblast/epiblast decision. Mesoderm before it splits into paraxial, lateral and intermediate parts; broad, largely unsupported~\cite{griffin1998,hug1997,schultemerker1994}. & \emph{ved, mespab, tbx16, eve1, her1, mespaa} \\
presomitic mesoderm & 30\% epiboly -- 6-somite & ZFIN: segmental plate; presumptive segmental plate; paraxial mesoderm; presumptive paraxial mesoderm & Paraxial mesoderm; from mesoderm & lateral margin, then segmental plate; Wnt, FGF, Notch oscillator & Paraxial mesoderm, the source of somites. Tbx6/tbx16 paraxial mesoderm; segmental plate with the her/Notch clock from bud on~\cite{hug1997,sawada2000,oates2002,holley2002}. & \emph{aplnra, gadd45ba, mespab, tbx16, msgn1, ripply1} \\
\rowcolor{gray!10}tailbud & bud -- 6-somite & ZFIN: tail bud & Tailbud; from margin & ventral margin converging at the vegetal pole; Wnt, FGF, BMP & Posterior progenitors (neuromesodermal) at the end of gastrulation. Tailbud with cdx and ntl/tbxta genes; feeds the posterior body~\cite{martin2010,joly1992,shimizu2005}. & \emph{vox, hes6, p4ha2, foxd5, znfl2a, nfkbiaa} \\
adaxial somite & bud -- 6-somite & ZFIN: somite; adaxial cell; myotome; fast muscle cell & Paraxial mesoderm (adaxial cells, somites); from presomitic mesoderm & next to the notochord; Hedgehog from the axis & First muscle and somite fates. Adaxial cells (slow muscle precursors) and forming somites~\cite{weinberg1996,coutelle2001,shanmugalingam1998,wohlgemuth2007}. & \emph{ved, ripply1, mespaa, gadd45bb, cdkn1ca, rdh10a} \\
\rowcolor{gray!10}lateral plate & bud -- 6-somite & ZFIN: lateral plate mesoderm; anterior lateral plate mesoderm; posterior lateral plate mesoderm; lateral mesoderm (+2 more) & Lateral plate mesoderm; from mesoderm & ventral and lateral margin; BMP & Lateral plate decisions (blood, vessels, heart, kidney) begin at bud. Lateral plate mesoderm with haemato-vascular precursors~\cite{gering1998,brown2000,mudumana2008}. & \emph{gata6, aplnrb, lmo2, blf, gata2a, tbx1} \\
intermediate mesoderm & bud -- 6-somite & ZFIN: pronephric mesoderm; presumptive pronephric mesoderm; intermediate mesoderm & Lateral plate mesoderm (kidney precursors); from mesoderm & between paraxial and lateral plate mesoderm & Kidney precursors, a lateral fate. Intermediate (pronephric) mesoderm~\cite{mudumana2008,pfeffer1998}. & \emph{ftr83, osr1, ftr82, acy3.2, lhx1a, twist1a} \\
\rowcolor{gray!10}heart & 6-somite & ZFIN: heart primordium & Lateral plate mesoderm (heart precursors); from lateral plate & anterior lateral plate mesoderm & Heart precursors at the end of the range (6-somite). Heart field (nkx2.5/2.7, hand2)~\cite{lee1996,yelon2000,novikov2013}. & \emph{nkx2.7, tmem88a, tbx20, nkx2.5, rbfox1l} \\
\midrule
\multicolumn{7}{@{}l}{\textit{axial}} \\
prechordal plate & 50\% epiboly -- 6-somite & ZFIN: prechordal plate; anterior axial hypoblast & Prechordal plate; from dorsal organizer & first internalized dorsal cells; highest Nodal & The anterior axial fate, which patterns the forebrain. Prechordal plate (gsc, frzb); its anterior part becomes the hatching gland~\cite{gritsman2000,thisse1995,seiliez2006}. & \emph{gsc, tagln2, pitx2, elovl6, icn, ctslb} \\
\rowcolor{gray!10}hatching gland & bud -- 6-somite & ZFIN: polster; hatching gland & Prechordal plate (anterior part); from prechordal plate & anterior prechordal plate (polster) & A terminal fate of the prechordal plate. Hatching gland (he1a, ctsl): a clear switch-like program from bud~\cite{vogel1997}. & \emph{gsc, klf17, pfn1, tagln2, pitx2, cldn7b} \\
axial mesoderm & 50\% epiboly -- 6-somite & ZFIN: axial chorda mesoderm; axial mesoderm; axial hypoblast & Notochord (axial mesoderm); from dorsal organizer & dorsal midline; Nodal, then low Wnt & The posterior axial fate before the notochord differentiates. Axial chorda mesoderm (noto/flh, foxa2) during gastrulation~\cite{strahle1993,talbot1995,schultemerker1994}. & \emph{noto, ta, serpinh1b, tob1a, cnn2, perp} \\
\rowcolor{gray!10}notochord & bud -- 6-somite & ZFIN: notochord; notochord posterior region & Notochord; from axial mesoderm & dorsal midline, under the neural plate & Differentiated axial structure; source of Hedgehog. Notochord (col2a1, col8a1, ntd5) from bud on~\cite{talbot1995,yan1995,gansner2008}. & \emph{ta, xbp1, zgc:162964, serpinh1b, cnn2, waslb} \\
\midrule
\multicolumn{7}{@{}l}{\textit{ectoderm}} \\
neural plate & 50\% epiboly -- 6-somite & ZFIN: neural plate; presumptive neural plate; anterior presumptive neural plate; neuroectoderm (+4 more) & Neural ectoderm; from blastomere & dorsal ectoderm; low BMP; recorded in ZFIN from 50\% epiboly & The neural/non-neural decision of the epiblast, then the neural plate before and during A-P subdivision. Presumptive neural plate, neuroectoderm, neural plate, keel and rod; broad, mostly unsupported (covers all neural regions)~\cite{dee2008,rimini1999,grinblat1998}. & \emph{fzd8b, her3, her6, hspb1, irx7, rx3} \\
\rowcolor{gray!10}forebrain eye & 75\% epiboly -- 6-somite & ZFIN: anterior neural plate; anterior neural keel; anterior neural rod; forebrain neural plate (+8 more) & Neural ectoderm (anterior); from neural plate & anterior neural plate; Wnt off & A-P subdivision of the neural plate: anterior. Forebrain and eye field (six3, rx, hesx1/anf, lhx2)~\cite{chuang1999,seo1998a,kazanskaya1997,ando2005}. & \emph{fzd8b, stm, six3b, rx3, hesx1, zic2a} \\
midbrain hindbrain boundary & 75\% epiboly -- 6-somite & ZFIN: midbrain hindbrain boundary neural plate; midbrain hindbrain boundary neural keel; midbrain hindbrain boundary neural rod; presumptive midbrain hindbrain boundary (+4 more) & Neural ectoderm (midbrain-hindbrain boundary); from neural plate & midbrain-hindbrain boundary; FGF8 organizer & A-P subdivision: the isthmic organizer. Midbrain-hindbrain boundary (eng, pax2/5/8, gbx2)~\cite{ekker1992,pfeffer1998,kikuta2003,reifers1998}. & \emph{her3, ism1, her5, mdkb, her11, shisa2} \\
\rowcolor{gray!10}hindbrain & 75\% epiboly -- 6-somite & ZFIN: hindbrain neural plate; hindbrain neural keel; hindbrain neural rod; presumptive hindbrain (+8 more) & Neural ectoderm (hindbrain); from neural plate & posterior neural plate; FGF, RA & A-P subdivision: hindbrain rhombomeres. Hindbrain (krox20/egr2b, hox genes)~\cite{oxtoby1993,moens1998,prince1998}. & \emph{irx7, her2, cyp26c1, efnb2a, dbx1a, hoxb1b} \\
primary neurons & bud -- 6-somite & ZFIN: primary neuron; primary neuron hindbrain; primary motor neuron; Rohon-Beard neuron (+6 more) & (extends Table 1: neurogenesis in the neural ectoderm); from neural plate & proneural clusters; Notch lateral inhibition & The first neuronal differentiation. Primary neurons (motor, Rohon-Beard, interneurons) from bud on~\cite{takke1999}. & \emph{pnocb, pnx, elavl3, neurod4, smox, isl1} \\
\rowcolor{gray!10}non neural ectoderm & 75\% epiboly -- 6-somite & ZFIN: non neural ectoderm; integument & Non-neural ectoderm; from blastomere & ventral epiblast; high BMP & The non-neural side of the ectoderm decision. Non-neural ectoderm, future epidermis (tp63, gata3); its early genes (dlx3b, foxi1) are pooled into the placodes~\cite{bakkers2002,neave1995,read1998}. & \emph{gata3, foxi1, tfap2a, tp63, foxi3a, tfap2c} \\
neural crest & 3-somite -- 6-somite & ZFIN: neural crest; neural crest cell; cranial neural crest; cranial neural crest cell & Neural plate border; from neural plate, non neural ectoderm & neural plate border; intermediate BMP, Wnt, FGF & The border decision: neural crest. Neural crest (sox10, tfap2a, foxd3) from 3-somite~\cite{dutton2001,knight2003,licornell2007,hernandezlagunas2005}. & \emph{tfap2a, foxd3, ovol1, cdh6, fam212aa, fam46ba} \\
\rowcolor{gray!10}placodes & 75\% epiboly -- 6-somite & ZFIN: preplacodal ectoderm; otic placode; trigeminal placode; neurogenic placode (+2 more) & Neural plate border; from neural plate, non neural ectoderm & neural plate border, anterior; preplacodal ectoderm & The border decision: placodes. Preplacodal ectoderm and placodes (six1, eya1, dlx3b, foxi1)~\cite{solomon2002,solomon2003,sahly1999,bessarab2004}. & \emph{foxi1, dlx3b, cldna, ptgs2a, tp63, pax8} \\
\midrule
\multicolumn{7}{@{}l}{\textit{signalling profiles}} \\
Nodal &  & readouts, ligands and antagonists (signalling table) &  & margin, then dorsal mesendoderm and axial mesoderm & General (received and sent in the same cells). Nodal readouts, ligands (ndr1, ndr2; autoregulated) and lefty (feedback, secreted)~\cite{rebagliati1998,feldman1998,thisse1999,feldman2002}. & \emph{lft1, lft2, ndr1, ndr2} \\
\rowcolor{gray!10}BMP &  & readouts, ligands and antagonists (signalling table) &  & ventral gastrula, then tailbud & General (received and sent in the same cells). BMP readouts and the ventral ligands (bmp2b, bmp4, bmp7a), which are BMP targets themselves; szl (feedback, secreted)~\cite{kishimoto1997,nguyen1998,dick2000}. & \emph{id1, id3, bambia, szl, msx1b, bmp2b} \\
BMP (blocking) &  & secreted antagonists, degradation (signalling table) &  & dorsal organizer (chd, nog1), paraxial mesoderm (fsta) & Where the signal is blocked (secreted antagonists). BMP antagonists made by cells with low BMP; shape the dorsal-ventral gradient~\cite{hammerschmidt1996,furthauer1999,yabe2003}. & \emph{chd, nog1, nog2, fsta} \\
\rowcolor{gray!10}BMP organizer (sending) &  & ligands, synthesis (signalling table) &  & dorsal organizer, then axial mesoderm & Where the signal is sent (ligands). The organizer's BMP ligand admp, which goes with low BMP readouts~\cite{lele2001,willot2002}. & \emph{admp} \\
Wnt &  & readouts, ligands and antagonists (signalling table) &  & margin, then tailbud & General (received and sent in the same cells). Wnt/beta-catenin readouts and the canonical ligands (wnt8a, wnt3a)~\cite{kelly1995,lekven2001,shimizu2005}. & \emph{sp5l, sp5a, notum1a, lef1, wnt8a, wnt3a} \\
\rowcolor{gray!10}Wnt (blocking) &  & secreted antagonists, degradation (signalling table) &  & organizer and prechordal plate, then anterior neural plate & Where the signal is blocked (secreted antagonists). Wnt antagonists (dkk1b, frzb, sfrp1a); protect the anterior from Wnt~\cite{hashimoto2000dkk1,tendeng2006}. & \emph{dkk1b, frzb, sfrp1a} \\
Wnt/PCP (sending) &  & ligands, synthesis (signalling table) &  & margin and mesoderm & Where the signal is sent (ligands). Non-canonical Wnt ligands (wnt5b, wnt11) of convergent extension; no readouts~\cite{heisenberg2000,kilian2003}. & \emph{wnt5b, wnt11} \\
\rowcolor{gray!10}FGF &  & readouts, ligands and antagonists (signalling table) &  & margin, then tailbud and midbrain-hindbrain boundary & General (received and sent in the same cells). FGF readouts and ligands (fgf8a, fgf24, fgf3, fgf4)~\cite{furthauer1997,reifers1998,draper2003}. & \emph{etv4, il17rd, fgf8a, fgf24, fgf3, fgf4} \\
Notch (receiving) &  & readouts (signalling table) &  & neural plate, presomitic mesoderm & Where the signal is received (readouts). Her/hey targets of Notch. & \emph{her4.1, her4.2, her4.3, her15.1, her15.2, her12} \\
\rowcolor{gray!10}Notch (sending) &  & ligands, synthesis (signalling table) &  & neural plate (primary neurons), presomitic mesoderm & Where the signal is sent (ligands). Delta ligands (dla, dlb, dlc, dld): lateral inhibition and the segmentation clock~\cite{haddon1998,holley2002}. & \emph{dla, dlb, dlc, dld} \\
RA (receiving) &  & readouts (signalling table) &  & central nervous system (somite stages) & Where the signal is received (readouts). Dhrs3a (feedback); a weak readout. & \emph{dhrs3a} \\
\rowcolor{gray!10}RA (sending) &  & ligands, synthesis (signalling table) &  & involuting and paraxial mesoderm & Where the signal is sent (synthesis). RA synthesis (aldh1a2)~\cite{begemann2001,grandel2002}. & \emph{aldh1a2} \\
RA (blocking) &  & secreted antagonists, degradation (signalling table) &  & anterior neural plate, margin, then tailbud & Where the signal is blocked (degradation). RA degradation (cyp26a1); RA-inducible, but its pattern follows anterior identity~\cite{kudoh2002,dobbsmcauliffe2004}. & \emph{cyp26a1} \\
\rowcolor{gray!10}Hedgehog (sending) &  & ligands, synthesis (signalling table) &  & axial mesoderm, notochord, floor plate & Where the signal is sent (ligands). Hedgehog ligands (shha, shhb); the readouts are not heterogeneous~\cite{krauss1993,ekker1995}. & \emph{shha, shhb} \\
\midrule
\multicolumn{7}{@{}l}{\textit{cell states}} \\
cell cycle s &  & literature table &  & all proliferating cells & A known non-identity axis, to be separated from identity. S phase genes~\cite{tirosh2016}. & \emph{pcna, mcm2, mcm5, mcm6, cdc6, slbp} \\
\rowcolor{gray!10}cell cycle g2m &  & literature table &  & all proliferating cells & A known non-identity axis, to be separated from identity. G2/M phase genes~\cite{tirosh2016}. & \emph{cdk1, ube2c, top2a, nusap1, cenpf, tpx2} \\
heat shock &  & literature table &  & handling of embryos and cells & Technical control (stress during handling). Heat shock response (hsp70, hsp90)~\cite{krone1994,lele1997,vihervaara2017}. & \emph{hsp70.3, hsp90aa1.2, bag3, dnajb1a, hsph1} \\
\rowcolor{gray!10}p53 response &  & literature table &  & stressed or damaged cells & Technical and biological stress control. P53 targets (mdm2, cdkn1a, ccng1)~\cite{fischer2017,langheinrich2002,robu2007}. & \emph{mdm2, cdkn1a, phlda3, tp53inp1, gadd45aa, rrm2b} \\
\midrule
\multicolumn{7}{@{}l}{\textit{controls}} \\
whole organism &  & ZFIN: whole organism only &  & whole embryo & Negative control: uniform genes should not vary between cells. Genes ZFIN records only in the whole organism from high to shield. & \emph{yes1, hdac1, ikzf5, fyna, stk11, dgcr8} \\
\rowcolor{gray!10}YSL & high -- 6-somite & ZFIN: yolk syncytial layer; external yolk syncytial layer; internal yolk syncytial layer (genes of curated programs left out) & YSL & syncytium between the blastoderm and the yolk; source of Nodal signals & Negative control: removed with the yolk before dissociation. No YSL cells are expected; genes of the curated programs and other controls are left out, so that only YSL-specific genes remain~\cite{rodaway1999}. & \emph{apoeb, tnfaip8l2b, tfa, sesn3, pleca, rgs4} \\
IFN &  & readouts (signalling table) &  & whole embryo & Negative control: should not vary between cells. Type I interferon response; no immune system yet. & \emph{isg15, mxa} \\
\rowcolor{gray!10}AhR &  & readouts (signalling table) &  & whole embryo & Negative control: should not vary between cells. Aryl hydrocarbon receptor (xenobiotic) response; embryo-wide. & \emph{cyp1a, cyp1b1} \\
GR &  & readouts (signalling table) &  & whole embryo & Negative control: should not vary between cells. Glucocorticoid receptor response; embryo-wide. & \emph{fkbp5, klf9, tsc22d3} \\
\rowcolor{gray!10}circadian &  & readouts (signalling table) &  & whole embryo & Negative control: should not vary between cells. Circadian clock; synchronised across the embryo. & \emph{per1b, per2, cry1aa, tefa} \\
dissociation &  & readouts (signalling table) &  & whole embryo & Technical control: stressed cells. Immediate early genes induced by tissue dissociation; flags stressed cells. & \emph{fosab, junba, jun, atf3, egr1} \\
\end{longtable}
